\section{Seed study: registered protocol}
\label{app:seed}

The two-model study that motivated the evaluation registered the procedures below; Appendix~\ref{app:seed-results} reports its results.

\paragraph{Data and split.}
We use NIH ChestX-ray14 \citep{wang2017chestxray}, whose findings are mined from radiology reports. A deterministic hash of patient identity assigns rows to disjoint train, validation, and test partitions. The registered manifest contains 18,212 training images and 5,343 test images from 1,659 test patients. The three cells pair LLaVA-1.5-7B with Effusion and Edema, and Qwen2.5-VL-7B with Effusion. All reported uncertainty clusters by patient. Prompts appear in Appendix~\ref{app:protocol}.

\paragraph{Models and consumed loci.}
For LLaVA-1.5-7B \citep{liu2024improved} and Qwen2.5-VL-7B \citep{bai2025qwen25vl}, we use the final visual block consumed by the connector or merger. Writing $\alpha$ for the signed intervention dose, synthetic hook tests establish downstream-path membership, a bitwise $\alpha=0$ no-op, and changed downstream representations and logits at nonzero dose. Probe and intervention therefore share one consumed representation.

\paragraph{Controlled decodability.}
For image $i$, let $\hvec_{it}\in\R^D$ denote the $D$-dimensional representation of visual token $t$ at the consumed final block; mean pooling over tokens yields $\hvec_i$. A fixed seed-0 Gaussian matrix $R\in\R^{D\times512}$ projects every model to a common readout dimension; featurewise scale $\boldsymbol{s}$ is fitted on training rows only. We then fit logistic regression with $C=1$, maximum 2,000 iterations, and seed 0. Each of 20 control seeds maps the recurring input type (AP view indicator, sex indicator, and age decade, giving 39 strata) to a random label. Let $y$ be the real test labels and $r$ the fitted probe's continuous scores. For control seed $k\in\{0,\ldots,19\}$, $y_k^{\mathrm{ctrl}}$ denotes the random labels and $r_k^{\mathrm{ctrl}}$ the corresponding control probe scores. Selectivity compares their AUROCs:
\begin{equation}
S = \operatorname{AUROC}(y,r)-\frac{1}{20}\sum_{k=0}^{19}\operatorname{AUROC}(y_k^{\mathrm{ctrl}},r_k^{\mathrm{ctrl}}).
\label{eq:selectivity}
\end{equation}
We report the real AUROC, control distribution, and $S$ with 2,000 patient-cluster bootstrap resamples. A positive interval for $S$ means that the target is more readable than same-capacity nuisance-type label maps. It does not yet imply downstream use. Paired contrasts reuse identical ordered test rows and bootstrap draws.

\paragraph{Probe-normal intervention.}
Let $\boldsymbol\beta_c$ be the projected-space logistic coefficient for the positive class of concept $c$. We lift it to the consumed $D$-dimensional representation and normalize it exactly as
\begin{equation}
\hat\wvec_c = \frac{R\left(\boldsymbol\beta_c\oslash\max(\boldsymbol{s},10^{-8})\right)}{\left\lVert R\left(\boldsymbol\beta_c\oslash\max(\boldsymbol{s},10^{-8})\right)\right\rVert_2}.
\label{eq:lift}
\end{equation}
The positive-class coefficient fixes the direction's clinical pole. At every visual token $t$, the registered intervention is
\begin{equation}
\hvec'_{it}=\hvec_{it}+\alpha\lVert\hvec_{it}\rVert_2\hat\wvec_c.
\label{eq:intervention}
\end{equation}
Here the sign of $\alpha$ selects the positive or negative clinical pole, while $|\alpha|$ is the perturbation norm as a fraction of each token's activation norm. At the first answer position, $P(\mathrm{yes})=\sigma(\ell_{\mathrm{yes}}-\ell_{\mathrm{no}})$, where each log probability is the maximum over case and leading-space variants that tokenize to one token; the endpoint is its paired mean change from the common unsteered baseline. Original cells use concept doses $\{-1,-0.5,-0.25,-0.1,0,0.1,0.25,0.5,1\}$ on 200 held-out test images. Controls are evaluated at the six nonzero doses excluding $\pm0.1$: 20 isotropic random unit directions, one sham formed by a fixed coordinate permutation of $\hat\wvec_c$ and renormalization, and fixed alternative clinical probe normals. Every control receives the same signed dose on the same rows. Dose monotonicity over $[-0.5,0.5]$ is descriptive.

\paragraph{Original-cell decision rule.}
Let $\Delta_c(\alpha)$ be the target direction's paired mean probability change, and let $\mathcal A=\{-1,-0.5,-0.25,0.25,0.5,1\}$ be the six controlled doses. The primary effect is $\max_{\alpha\in\mathcal A}\operatorname{sign}(\alpha)\Delta_c(\alpha)$, attained at dose $\alpha^\star$. A cell clears the generic intervention gate when this selected effect exceeds the 95th percentile of the 20 random directions' mean probability changes at that same $\alpha^\star$, each multiplied by $\operatorname{sign}(\alpha^\star)$, and the maximum absolute sham effect over $\mathcal A$. The random directions are not separately maximized over doses. This asymmetric selection defines a registered reference gate, not a cross-dose-corrected 5\% test. The conjunction is fixed across the three cells. The independent supplement below instead evaluates a dose locked before observing its patients.

\paragraph{Direction-specificity supplement.}
The initial experiment prespecified the random and sham comparisons; its clinical-direction analysis was descriptive. The Nodule response motivated a separately registered supplement at the locked Qwen dose, $\alpha=+0.25$. We select one row by label-blind hash for each of 400 ChestX-ray14 test patients absent from the original 200-image intervention cohort. The semantic family contains all five non-target normals from the pre-existing Qwen direction bundle (Atelectasis, Pneumothorax, Cardiomegaly, Mass, and Nodule), rather than directions chosen from the supplement outcome. Each uses the same training rows, projection, scaling, classifier, lifting rule, and normalization as Effusion. Together with a common baseline, 20 random directions, and the sham, the grid produces 28 configurations and 11,200 per-image outcomes; exact row selection appears in Appendix~\ref{app:seed-results}.

Semantic ownership is a contrast on a direction--question response surface, not a property of a direction in isolation. Let $W_{q,d}$ denote the paired mean change in question $q$'s yes probability after intervening with direction $d$ at the locked dose. A large $W_{q,d}$ establishes that $d$ can write $q$'s answer; it does not establish that $d$ owns that answer. Within a fixed clinical family, the ownership contrast is
\begin{equation}
O_q=W_{q,q}-\max_{d\neq q}W_{q,d}.
\label{eq:margin}
\end{equation}
The registered supplement estimates the Effusion row of this surface, so the reported ownership contrast $O_{\mathrm{Effusion}}$ is Effusion minus its maximum alternative. The maximum is recomputed inside each of 5,000 patient-bootstrap draws. Thus $O_{\mathrm{Effusion}}>0$ means Effusion beats every member of the fixed clinical family, while $O_{\mathrm{Effusion}}<0$ means at least one alternative is stronger. Adding competitors can only maintain or decrease $O_q$: a positive advantage is conditional on the candidate family, whereas a negative counterexample does not require an exhaustive family. This operational ownership criterion measures relative advantage, not a necessary condition for a concept to participate in computation. Direction specificity requires a positive Effusion effect above random p95 and absolute sham, together with a one-sided 95\% lower bound for $O_{\mathrm{Effusion}}$ above zero.

\paragraph{Causal-ownership matrix.}
To test whether another clinical normal better matches the answer pathway, we prospectively cross six Qwen directions (Effusion, Atelectasis, Pneumothorax, Cardiomegaly, Mass, and Nodule) with their six matching questions on 400 new patients. Every question uses the same $\alpha=+0.25$, 119 random directions, and a sham. For question $q$, ownership requires the diagonal effect $W_{q,q}$ to exceed every other clinical direction, every random direction, and the absolute sham effect. A centered studentized max-$T$ bootstrap gives simultaneous 95\% lower bounds over all 30 clinical contrasts. The shared-alias family contains two statistics per direction: mean excess over the matching diagonals on its five off-diagonal questions and signed mean effect on those questions. These give 12 max-$T$ bounds; the effect must also exceed the random alias maximum and maximum absolute sham effect, detailed in Appendix~\ref{app:seed-results}.

For a family of $m$ statistics, let $\widehat C_j$ be statistic $j$'s estimate and $C^*_{bj}$ its value in joint patient-bootstrap draw $b$, with $B=5{,}000$ draws. Write $s_j=\operatorname{SD}_{b}(C^*_{bj})$ for the bootstrap standard deviation with denominator $B-1$, held fixed across draws, and $Q_{.95}$ for the empirical 95th percentile with linear interpolation. For $s_j>10^{-12}$, the centered standardized draw $Z_{bj}$, common critical value $c_{.95}$, and simultaneous lower bound $L_j$ are
\begin{equation}
Z_{bj}=\frac{C^*_{bj}-\widehat C_j}{s_j},\qquad
c_{.95}=Q_{.95}\!\left(\left\{\max_{1\leq j\leq m}Z_{bj}\right\}_{b=1}^{B}\right),\qquad
L_j=\widehat C_j-c_{.95}s_j.
\label{eq:max-t}
\end{equation}
Joint resampling preserves paired conditions; no standard error is re-estimated within a draw. Appendix~\ref{app:simultaneous-bounds} specifies the families, degenerate components, and the separate percentile intervals for target-minus-maximum margins.

\paragraph{Answer-encoding crossover.}
We next register a paired sensitivity study on 200 patients from the causal-ownership matrix cohort; the other 200 form a patient-disjoint capability-calibration set. This is not the earlier Effusion specificity cohort. Six questions cross the original yes/no prompt with A-present and B-present prompts, which exchange the finding-to-letter mapping. Each uses both $\alpha=\pm0.25$, six clinical directions, 20 random directions, and a sham. A question enters the primary endpoint only if both clean coded mappings have AUROC one-sided 95\% lower bounds above 0.5 and at least ten patients per class on calibration. At $\alpha=+0.25$, we compare squared mean effects in components that retain or reverse raw A-minus-B sign under the mapping exchange, using an unbiased estimator and 2,000 joint patient-bootstrap draws. Appendix~\ref{app:seed-results} defines the components and full control rule. Descriptive clinical competition uses the contrast in Equation~\ref{eq:margin}, replacing probability change with finding-present logit-margin change within each prompt.

\paragraph{Independent Mass confirmation.}
We prospectively test Mass on 200 patients absent from all preceding Qwen intervention cohorts, retaining the crossover's 200 calibration patients. Two question wordings (\textit{show}, \textit{is}) cross explicit yes/no, A-present, and B-present instructions; the original yes/no prompt is a seventh anchor. At $\alpha=+0.25$, each cell receives six clinical directions, 119 random directions, and a sham. The four coded cells share 20 clinical comparisons with simultaneous 95\% lower bounds. Confirmation requires positive bounds and Mass effects above every random effect and absolute sham, either across all four coded cells or both original-wording cells. Appendix~\ref{app:seed-results} specifies calibration and paired descriptive contrasts.

